% !TeX root = ../papers/nc-p3.tex
\section{二十の中心指標での通常二重点と局所行列表示}
\label{sec:twenty-central-nodes}
\begin{theorem}\label{thm:twenty-central-nodes}
十五次元成分の一般窓の射影四次中心上で、前節の二十指標は全て通常二重点である。
その各点の単純加群の次元は4と12である。中心をその点で完備化し、
可逆な四次変数の値を1にして次数方向を除くと、中心と折り畳み環は
\[
 R\simeq k[[a,b,c,d]]/(ad-bc),\qquad
 B^{\wedge}\simeq\operatorname{End}_R(R^{\oplus4}\oplus I^{\oplus12})
\]
を満たす。$I$は一般二行二列行列の余核である（座標の選び方による）。
特にこの二十点はそれぞれ孤立している。
\end{theorem}
\begin{proof}
$A=e_0Be_0$とし、長さ4のループの固定した整数線型結合$u\in A$を取る。
その中心分数体上での四次被約特性多項式を$H(T)$とする。
中心は正規で、整な元の固有値の基本対称式も整であるから$H\in Z[T]$である。
四次中心の第五変数$z$の$P$上のモニック関係を$F$とする。$A$の$P$自由階数は64で、
中心四次拡大の一般分裂体では四つの四行行列環の直積になるので、補助変数$v$を入れて
\[
 \det\nolimits_{P}(T+vL_z-L_u)
   =\operatorname{Norm}_{Z/P}\bigl(H(T+vz)\bigr)^4,\qquad
 \det\nolimits_P(T-L_z)=F(T)^{16}.
\]
$L$は左正則作用である。これらは係数環上の多項式恒等式で、重根を持つ有限中心
ファイバーにも特殊化できる。特に中心根$\alpha$が$F$で重複度$e$なら、
$H_\alpha$の根$\lambda$の重複度$m$は、任意の定数$v$に対して
$L_u-\lambda+v(L_z-\alpha)$の零固有値に少なくとも$4em$を寄与する。
これはArtin中心のノルムが各剰余値を局所長だけ掛けた積であることによる。

この恒等式を有理点$t$で$p=2657$へ還元する手続きを明記する。
既定の四中心変数と$z$は局所整数環上に取り、四変数を零とした商で
規定の剰余基底と$1,z,z^2,z^3$の独立性を検算する。長さ13の剰余商は零である
（前節までのフィルター商の、零パラメータでの長さ13と長さ9の同一次元および
安定埋め込みによる。斉次剰余商の検算も下記コードに含める）。
従ってこれらの経路リフトから$P$自由加群の全射を得る。標数零で既知の
Hilbert関数に一致し、自由加群の部分加群である核はtorsion-freeなので、
局所整数環でも同型である。中心四元を含む斉次剰余基底へ延長でき、
$z^4$の関係も整係数となる。この整モデル$P[z]/F$は正規である。
実際、標数零では既知の正規環であり、特殊ファイバーも$(1,2,3,4)$で
$F$が四つの単根を持つので一般に被約である。従って垂直な高さ1素点では
$p$が極大イデアルを生成し、水平な高さ1素点でも正則である。
超曲面なので$S_2$であり、平坦性と標数零での整性から整域でもある。
従って$H$の係数はこの整モデルに入り、上の恒等式を剰余へ移せる。

二十の番号付き四組に四次中心元と$u$を評価し、$x_1=1$に規格化する。
全四組で$x_1$の剰余は非零である。$z,u$の値を$\alpha,\lambda$とすると、
第0隅の64次元正則作用で$L_z-\alpha$の零固有値の重複度は各々32となる。
$v=1,3,7$の三値で$L_u-\lambda+v(L_z-\alpha)$の零固有値重複度の
最小は各々8である。従って$e=2$であり、上の下界から$m\le1$となる。
$\lambda$は$H_\alpha$の根である。実際、$H(u)=0$を一次元の隅表現で評価すればよい。
従って$H_\alpha'(\lambda)\ne0$かつ$F''(\alpha)\ne0$である。
\path{computations/verify_central_nodes.py}は$u$の係数と以上の計算を再現する。

この二つの非消滅は各番号について標数零の$t$に持ち上がり、成分の一般窓へ移る。
ここで$H$を族でも正則に取れることを補足する。$t$の近くの滑らかなテンソル基底上で、
$P$自由表示、$Z=P[z]/F$、有限な経路の乗法表を前節までの単元小行列で取る。
全中心族は平坦で幾何学的に正規なファイバーを持ち、基底が滑らかなので正規である。
一般商の被約多項式の係数は再び整性からその環に入る。各ファイバーの一般の四変数の値で
$z$の四根は相異なり、上のノルム恒等式で$v$を不定として四根ごとの因子を分離すれば、
特殊化した$H$はその隅の被約多項式に一致する。実際、各因子の$T$の根は
$v$について傾き$-\alpha$を持ち、同一の傾きのモニック因子の積は一意である。
これで特異な中心指標にも適用できる。同様に、$H'_\alpha(\lambda)\ne0$の下で
最小単純に対応する原始冪等元の一般行列階数は1となる。詳しくは、
指標での被約traceは有限ファイバーの根基上で消える。根基元の高い冪は
中心極大イデアルと$A$の積に入り、そこでは全ての正冪のtraceが剰余で零となる。
Newton恒等式により根基元の全固有値は零である。従って誘導traceは各単純行列因子で
通常のtraceの定数倍となり、原始冪等元における値は持ち上げの一般行列階数である。
再びNewton恒等式から$H_\alpha$は各単純の特性多項式をこの階数だけ掛けた積となる。
$S e_0$は一次元で$u$の値が$\lambda$だから、その重複度1により前節の
$(1,2),(2,1)$の可能性が消え、残る単純の全次元は12となる。

次に$x_1=1$とし、$R$をこの射影中心の局所完備環とする。
点は特異であり、$F''(\alpha)\ne0$より四変数正則局所環中の超曲面の
重複度はちょうど2である。完備化した折り畳み環の二つの単純に対応する
原始冪等元を$p,q$と書く。ともに一般行列階数1なので一般商は分裂し、
$pB^{\wedge}p=R$となる。実際これは分数体の中で$R$上整な有限環である。
$B^{\wedge}$は$R$上極大Cohen--Macaulayなので反射的である。
この性質と高さ1でのAzumaya性は前節のエタールな次数自明化と完備化にも保たれる。
より詳しくは、完成写像の形式ファイバーが幾何学的に正則なので、完成した中心の
高さ1の点ではその下の点も正則である。よって左右乗法から
\[
 B^{\wedge}=\operatorname{End}_R(B^{\wedge}p)
\]
を得る。一般商では同型であり、高さ1のAzumaya環では階数1の$p$が充満で
森田同型となり、両辺の反射性から全体でも同型である。
$I=qB^{\wedge}p$と置けば原始冪等元の個数から
$B^{\wedge}p=R^{\oplus4}\oplus I^{\oplus12}$となる。
$I$は階数1の極大Cohen--Macaulay加群で、自由なら$p,q$が同値になり矛盾する。

$J=I^* I\subset R$を評価の像とする。基本代数
$\operatorname{End}_R(R\oplus I)$の$R$隅への射影は、$J$を割ると
$R/J$への代数全射になる。極大イデアルを$\mathfrak m$とすると$J=\mathfrak m$である。
実際$J$に単元があれば$I$は自由となる。一方$R/J$が非自明な余接空間を持てば、
その双対から剰余単純の非自明な自己一次拡大を得る。全射と森田同値で引き戻しても
非分裂であり、前節の$x_1=1$での最小単純の剛性に反する。中山より結論が従う。

最後に$R=S/(f)$、$S$を四変数の完備正則局所環とする。
$I$は$S$上深さ3なので最小自由表示は長さ1で、正方行列$M$による。
その行列式は単元倍の$f$である。実際、その素因子は加群の台から$f$だけであり、
$(f)$で局所化した長さは階数1ゆえ1である。$M$の成分は全て極大イデアルに入り、
$f$の重複度が2なのでサイズは高々2で、非自由性から丁度2である。
転置の核は$R$上で随伴行列の転置の像に一致する
（$S$に持ち上げ、積が$f$倍なら随伴行列を掛けて$f$を消去する）。
従って$I$の二生成元と双対との評価イデアル$J$は随伴行列、従って$M$の
四成分で生成される。$J=\mathfrak m$なのでこの四成分は四次元余接空間の基底をなし、
$S$の正則パラメータ系である。これを$a,b,c,d$とすれば主張の表示を得る。
\end{proof}

\begin{corollary}
一般の射影中心に残る特異点があれば、各点の単純加群は全て次元8で、
一種類または二種類である。二種類ならその完備環上の代数は
$\operatorname{End}_R(R^{\oplus8}\oplus J^{\oplus8})$であり、
$J$は非自由な階数1反射的加群である。
\end{corollary}
\begin{proof}
前節の完備化での階数和$4=\sum d_jm_j$を使う。
$d=4$があれば一種類だけで$m=1$である。全原始冪等元が同値で、
一般行列階数1の隅は正規中心に等しいため、完備化は十六次行列環である。
忠実平坦な降下によりAzumaya点となり、特異点に矛盾する。
$d=1$は既に二十点で尽き、$d=3$には$d=1$の相方が必要なので、
残る全ての$d_j$は2である。二種類なら両$m_j=1$であり、先の高さ1の
森田同型を反射性で延長する証明をそのまま適用できる。
\end{proof}
